\section{Estado actual}

\subsection{Volúmenes actuales}

\begin{center}
\small
\begin{tabular}{@{}p{54mm}p{28mm}p{56mm}@{}}
\toprule
\textbf{Relación} & \textbf{Filas} & \textbf{Nota} \\
\midrule
\cd{monitoreo\_sc\_electrico} & 36\,469 & 285 CSV, ventanas de 5\,min. \\
\cd{radiacion\_sc\_15s} & 94\,868 & Ventanas de 15\,s. \\
\cd{radiacion\_sc\_clearsky} & 94\,868 & Un cielo despejado por instante. \\
\cd{radiacion\_sc\_poa} & 56\,450 & Instantes con POA por arreglo. \\
\cd{lecturas\_ambientales\_sc} & $\approx$\,812\,000 & Ingesta de AgroDash, formato largo. \\
\bottomrule
\end{tabular}
\end{center}

La verificación cruzada confirmó que el crudo sigue siendo crudo (temperaturas en 85 y el
\emph{offset} presentes en la tabla base) y que las vistas hacen su trabajo (85 pasa a
NULL, la potencia máxima cae a un valor físico, la irradiancia negativa queda en cero,
kt* con percentil 95 en 1{,}01 y 99{,}4\% de filas aprobando el control de calidad).

\subsection{Decisiones que conviene no revisitar}

\begin{itemize}[leftmargin=*,itemsep=3pt]
  \item El crudo se guarda; la corrección vive en vistas. Cualquier propuesta de
    ``limpiar al insertar'' ya se descartó explícitamente.
  \item No se generan datos sintéticos para los huecos largos de 126 y 71 días. Como
    referencia paralela se usa NASA POWER.
  \item Los duplicados exactos por md5 se eliminan; los fragmentos numerados requieren
    decisión humana.
  \item Las bases de las dos regiones no se fusionan a nivel de almacenamiento, aunque un
    agente lea ambas.
\end{itemize}

\vspace{1em}
\begin{clave}[title=Dónde seguir]
El sistema de memoria del proyecto (\cd{docs/memoria/INDEX.md}) es la fuente viva: un
tema por archivo, con la evidencia contada y la fecha de cada decisión. Este documento
resume la arquitectura; ese índice explica \emph{por qué} cada pieza quedó así.
\end{clave}
